\chapter{Cod sursă relevant}
\label{ch:anexe}

În această anexă sunt incluse fragmente de cod din \texttt{WalletServiceImpl} și din suita de teste, la care se face referire pe parcursul documentației. Codul complet se găsește în repository-ul GitHub al proiectului.

\section{Fragmente din \texttt{WalletServiceImpl}}

\subsection{Constructorul cu injecție de dependențe}

\begin{lstlisting}[language=Java]
@Service
public class WalletServiceImpl implements WalletService {

    private static final int WARNING_THRESHOLD_PERCENT  = 75;
    private static final int CRITICAL_THRESHOLD_PERCENT = 90;
    private static final Set<String> ALLOWED_CURRENCIES =
        Set.of("RON", "EUR");

    private final TripRepository tripRepository;
    private final TripWalletRepository tripWalletRepository;
    private final UserRepository userRepository;
    private final WalletTransactionRepository walletTransactionRepository;

    public WalletServiceImpl(TripRepository tripRepository,
                             TripWalletRepository tripWalletRepository,
                             UserRepository userRepository,
                             WalletTransactionRepository walletTransactionRepository) {
        this.tripRepository = tripRepository;
        this.tripWalletRepository = tripWalletRepository;
        this.userRepository = userRepository;
        this.walletTransactionRepository = walletTransactionRepository;
    }
    // ... metodele clasei
}
\end{lstlisting}

\subsection{Metoda \texttt{addExpenseOwnedByUser}}

\begin{lstlisting}[language=Java]
@Override
@Transactional
public WalletTransaction addExpenseOwnedByUser(
        Long tripId, String usernameOrEmail,
        BigDecimal amount, WalletCategory category,
        String title, String note,
        LocalDate spentAt, Integer dayIndex) {

    TripWallet wallet = getOrCreateWalletOwnedByUser(tripId, usernameOrEmail);

    if (amount == null || amount.compareTo(BigDecimal.ZERO) <= 0) {
        throw new IllegalArgumentException("Suma trebuie sa fie > 0.");
    }
    String safeTitle = (title == null) ? "" : title.trim();
    if (safeTitle.isBlank()) {
        throw new IllegalArgumentException("Titlul este obligatoriu.");
    }
    if (dayIndex != null && dayIndex <= 0) {
        throw new IllegalArgumentException("Ziua trebuie sa fie >= 1.");
    }

    WalletTransaction tx = new WalletTransaction();
    tx.setWallet(wallet);
    tx.setAmount(amount);
    tx.setCategory(category != null ? category : WalletCategory.OTHER);
    tx.setTitle(safeTitle);
    tx.setNote(note);
    tx.setSpentAt(spentAt != null ? spentAt : LocalDate.now());
    tx.setDayIndex(dayIndex);

    return walletTransactionRepository.save(tx);
}
\end{lstlisting}

\subsection{Metoda \texttt{computeDaysRemainingSafe}}

\begin{lstlisting}[language=Java]
private Integer computeDaysRemainingSafe(Trip trip) {
    LocalDate start = trip.getStartDate();
    LocalDate end = trip.getEndDate();

    if (start == null || end == null || end.isBefore(start)) {
        return null;
    }

    LocalDate today = LocalDate.now();
    if (today.isAfter(end)) {
        return 0;
    }
    if (today.isBefore(start)) {
        return (int) (ChronoUnit.DAYS.between(start, end) + 1);
    }
    return (int) (ChronoUnit.DAYS.between(today, end) + 1);
}
\end{lstlisting}

\section{Fragmente din suita de teste}

\subsection{Configurare comună}

\begin{lstlisting}[language=Java, caption={Schelet din \texttt{...EquivalenceTest.java}}]
@ExtendWith(MockitoExtension.class)
@MockitoSettings(strictness = Strictness.LENIENT)
public class WalletServiceImplEquivalenceTest {

    @Mock private TripRepository tripRepo;
    @Mock private TripWalletRepository walletRepo;
    @Mock private UserRepository userRepo;
    @Mock private WalletTransactionRepository txRepo;

    @InjectMocks private WalletServiceImpl service;

    private TripWallet wallet;
    private Trip trip;
    private final Long tripId = 10L;
    private final String login = "user@test.com";

    @BeforeEach
    public void setUp() {
        UserEntity user = new UserEntity();
        user.setId(1L);
        user.setEmail(login);

        trip = new Trip();
        trip.setId(tripId);
        trip.setUser(user);
        trip.setStartDate(LocalDate.of(2026, 4, 15));
        trip.setEndDate(LocalDate.of(2026, 4, 22));

        wallet = new TripWallet();
        wallet.setId(100L);
        wallet.setTrip(trip);
        wallet.setBudgetTotal(new BigDecimal("1000"));
        wallet.setCurrency("RON");

        when(userRepo.findByEmail(login)).thenReturn(user);
        when(tripRepo.findById(tripId)).thenReturn(Optional.of(trip));
        when(walletRepo.findByTrip_Id(tripId))
            .thenReturn(Optional.of(wallet));
    }
}
\end{lstlisting}

\subsection{Test pentru clasa de echivalență P4 (sumă invalidă)}

\begin{lstlisting}[language=Java]
@Test
public void testAdaugaCheltuiala_SumaNegativa_AruncaExceptie() {
    BigDecimal sumaNegativa = new BigDecimal("-10");
    assertThrows(IllegalArgumentException.class, () ->
        service.addExpenseOwnedByUser(tripId, login,
            sumaNegativa, WalletCategory.FOOD,
            "Pranz", null, LocalDate.now(), 1));

    verify(txRepo, never()).save(any());
}
\end{lstlisting}

\subsection{Test pentru BVA pe pragul de 75\%}

\begin{lstlisting}[language=Java]
@Test
public void testSumar_ExactPragAtentie_75LaSuta() {
    wallet.setBudgetTotal(new BigDecimal("100"));
    when(txRepo.sumAmountByWalletId(any()))
        .thenReturn(new BigDecimal("75"));

    WalletSummary summary =
        service.computeSummaryOwnedByUser(tripId, login);

    assertEquals(75, summary.getPercent());
    assertEquals(WalletSummary.Status.ATENTIE, summary.getStatus());
    assertEquals(new BigDecimal("25.00"),
        summary.getRemaining().setScale(2));
}
\end{lstlisting}

\subsection{Test pentru Condition Coverage cu 3 operanzi OR}

\begin{lstlisting}[language=Java]
@Test
public void testStergere_OwnershipCheck_ConditionCoverage() {
    Long txId = 42L;

    // Caz 1: tranzactie cu wallet null -> primul operand TRUE
    WalletTransaction txCuWalletNull = new WalletTransaction();
    txCuWalletNull.setId(txId);
    txCuWalletNull.setWallet(null);
    when(txRepo.findById(txId)).thenReturn(Optional.of(txCuWalletNull));
    assertThrows(IllegalStateException.class, () ->
        service.deleteTransactionOwnedByUser(tripId, login, txId));

    // Caz 2: tranzactie cu id null -> al doilea operand TRUE
    WalletTransaction txCuIdNull = new WalletTransaction();
    txCuIdNull.setWallet(wallet);
    when(txRepo.findById(txId)).thenReturn(Optional.of(txCuIdNull));
    assertThrows(IllegalStateException.class, () ->
        service.deleteTransactionOwnedByUser(tripId, login, txId));

    // Caz 3: id-uri diferite -> al treilea operand TRUE
    TripWallet altWallet = new TripWallet();
    altWallet.setId(999L);
    WalletTransaction txCuIdGresit = new WalletTransaction();
    txCuIdGresit.setId(txId);
    txCuIdGresit.setWallet(altWallet);
    when(txRepo.findById(txId)).thenReturn(Optional.of(txCuIdGresit));
    assertThrows(IllegalStateException.class, () ->
        service.deleteTransactionOwnedByUser(tripId, login, txId));

    // Caz 4: toate operandele FALSE -> stergere reusita
    WalletTransaction txValid = new WalletTransaction();
    txValid.setId(txId);
    txValid.setWallet(wallet);
    when(txRepo.findById(txId)).thenReturn(Optional.of(txValid));
    service.deleteTransactionOwnedByUser(tripId, login, txId);
    verify(txRepo).delete(txValid);
}
\end{lstlisting}

\subsection{Test țintit pentru mutantul \textsc{MathMutator}}

\begin{lstlisting}[language=Java]
@Test
public void should_kill_percent_math_mutations() {
    wallet.setBudgetTotal(new BigDecimal("200"));
    when(txRepo.sumAmountByWalletId(any()))
        .thenReturn(new BigDecimal("50"));

    WalletSummary s = service.computeSummaryOwnedByUser(tripId, login);

    // 50 / 200 * 100 = 25%. Daca multiply si divide
    // sunt schimbate intre ele, calculul va da o valoare
    // diferita (in particular, 0 dupa rotunjire).
    assertEquals(25, s.getPercent());
}
\end{lstlisting}

\subsection{Test pentru Circuits pe \texttt{computeDaysRemainingSafe}}

\begin{lstlisting}[language=Java]
@Test
public void testCircuiteIndependente_computeDaysRemainingSafe() {
    // C1: start = null
    trip.setStartDate(null);
    trip.setEndDate(LocalDate.now().plusDays(5));
    WalletInsights i1 = service.computeInsightsOwnedByUser(tripId, login);
    assertNull(i1.getDaysRemaining());

    // C2: end = null
    trip.setStartDate(LocalDate.now());
    trip.setEndDate(null);
    WalletInsights i2 = service.computeInsightsOwnedByUser(tripId, login);
    assertNull(i2.getDaysRemaining());

    // C3: end < start (interval invalid)
    trip.setStartDate(LocalDate.now().plusDays(10));
    trip.setEndDate(LocalDate.now().plusDays(5));
    WalletInsights i3 = service.computeInsightsOwnedByUser(tripId, login);
    assertNull(i3.getDaysRemaining());

    // C4: today > end (calatorie incheiata)
    trip.setStartDate(LocalDate.now().minusDays(20));
    trip.setEndDate(LocalDate.now().minusDays(10));
    WalletInsights i4 = service.computeInsightsOwnedByUser(tripId, login);
    assertEquals(0, i4.getDaysRemaining());

    // C5: today < start (calatorie viitoare)
    trip.setStartDate(LocalDate.now().plusDays(3));
    trip.setEndDate(LocalDate.now().plusDays(9));
    WalletInsights i5 = service.computeInsightsOwnedByUser(tripId, login);
    assertEquals(7, i5.getDaysRemaining());

    // C6: caz normal (calatorie in desfasurare)
    trip.setStartDate(LocalDate.now().minusDays(2));
    trip.setEndDate(LocalDate.now().plusDays(4));
    WalletInsights i6 = service.computeInsightsOwnedByUser(tripId, login);
    assertEquals(5, i6.getDaysRemaining());
}
\end{lstlisting}

\section{CFG pentru \texttt{addExpenseOwnedByUser}}

Pentru completitudine, mai jos e prezentat și CFG-ul pentru \texttt{addExpenseOwnedByUser} (figura \ref{fig:cfg-addExpense}), care e cea mai testată metodă din clasa noastră.

\begin{figure}[!htb]
    \centering
    \includegraphics[width=0.7\linewidth]{03-cfg-addExpense.png}
    \caption{CFG pentru \texttt{addExpenseOwnedByUser}: trei decizii principale (sumă, titlu, dayIndex) cu trei branch-uri de eroare și o cale de succes.}
    \label{fig:cfg-addExpense}
\end{figure}
